\chapter{Testing and Results}

\section{Testing Strategy}

The testing approach for UPI-Guard++ covers three distinct areas: unit-level testing of the feature engineering functions, integration testing of the API endpoints, and model performance evaluation on held-out data.

\subsection{Unit Testing Feature Engineering}

We manually verified the correctness of the \texttt{build\_features()} function by creating small synthetic DataFrames with known values and checking that the output features matched expectations. This was particularly important for the velocity feature and the Z-score computation.

\begin{table}[H]
\centering
\caption{Feature Engineering Unit Test Cases}
\label{tab:unit_tests}
\begin{tabular}{|p{3.5cm}|p{3.5cm}|p{3.5cm}|c|}
\hline
\textbf{Test Case} & \textbf{Input} & \textbf{Expected Output} & \textbf{Pass?}\\
\hline
Night flag: 2 AM & hour = 2 & is\_night = 1 & \checkmark\\
\hline
Night flag: 10 AM & hour = 10 & is\_night = 0 & \checkmark\\
\hline
Salary week: day 3 & day = 3 & is\_salary\_week = 1 & \checkmark\\
\hline
Salary week: day 15 & day = 15 & is\_salary\_week = 0 & \checkmark\\
\hline
Cross-state: same & MH $\to$ MH & cross\_state = 0 & \checkmark\\
\hline
Cross-state: different & MH $\to$ DL & cross\_state = 1 & \checkmark\\
\hline
Z-score: mean amount & amount = mean & z $\approx$ 0 & \checkmark\\
\hline
Z-score: 3x mean & amount = 3$\times$mean & z $\approx$ 2 & \checkmark\\
\hline
Velocity: 5 txns in 1h & 5 rapid transactions & velocity = 5 & \checkmark\\
\hline
\end{tabular}
\end{table}

\subsection{Integration Testing API Endpoint}

The FastAPI endpoint was tested using a combination of the built-in Swagger UI (\texttt{/docs}) and direct Python requests. We created a set of test cases representing different risk profiles and verified that the API returned sensible predictions.

\begin{table}[H]
\centering
\caption{API Integration Test Cases}
\label{tab:api_tests}
\begin{tabular}{|p{2.5cm}|p{4cm}|l|l|c|}
\hline
\textbf{Test Case} & \textbf{Key Characteristics} & \textbf{Expected} & \textbf{Actual} & \textbf{Pass?}\\
\hline
Low-risk normal & Rs.500, daytime, established account, P2M, same state & Safe & Safe & \checkmark\\
\hline
High-risk mule & Rs.48000, 2 AM, 15-day account, 6 txns/hour, P2P & Fraud & Fraud & \checkmark\\
\hline
Cross-state large & Rs.25000, Maharashtra to Delhi, P2P & Fraud & Fraud & \checkmark\\
\hline
Moderate risk & Rs.5000, evening, moderate velocity & Safe & Safe & \checkmark\\
\hline
New account large & Rs.40000, account\_age=10 days & Fraud & Fraud & \checkmark\\
\hline
Regular grocery & Rs.800, P2M, Grocery, known state & Safe & Safe & \checkmark\\
\hline
Weekend large P2P & Rs.30000, Sunday, cross-bank & Fraud & Fraud & \checkmark\\
\hline
\end{tabular}
\end{table}

All seven integration test cases produced the expected classification, giving us confidence that the full pipeline from API request to prediction output is working correctly.

\section{Model Performance Results}

\subsection{Final Test Set Evaluation}

After training on the 64\% training split and selecting the threshold on the 16\% validation split, the model was evaluated once on the 20\% held-out test set. The results are presented below.

\begin{table}[H]
\centering
\caption{Final Model Performance Metrics on Test Set}
\label{tab:metrics}
\begin{tabular}{|l|r|}
\hline
\textbf{Metric} & \textbf{Score}\\
\hline
ROC-AUC & 0.9714\\
\hline
PR-AUC & 0.9168\\
\hline
Accuracy & 0.9752\\
\hline
Precision (Fraud class) & 0.8921\\
\hline
Recall (Fraud class) & 0.9138\\
\hline
F1-Score (Fraud class) & 0.9028\\
\hline
Optimal Threshold & 0.18\\
\hline
\end{tabular}
\end{table}

The ROC-AUC of 0.9714 indicates that the model correctly ranks a randomly chosen fraud transaction above a randomly chosen legitimate transaction about 97\% of the time. The PR-AUC of 0.9168 is particularly significant because it is robust to class imbalance and reflects genuine performance on the minority (fraud) class.

The optimal threshold of 0.18 (rather than the default 0.5) reflects the cost asymmetry: since missing a fraud case costs 25 times more than a false alarm ($C_{FN}/C_{FP} = 5000/200 = 25$), the model is tuned to be more aggressive in flagging suspicious transactions.

\subsection{Confusion Matrix}

\begin{table}[H]
\centering
\caption{Confusion Matrix on Test Set (approx. 20,000 records)}
\label{tab:confusion}
\begin{tabular}{|l|c|c|}
\hline
 & \textbf{Predicted: Safe} & \textbf{Predicted: Fraud}\\
\hline
\textbf{Actual: Safe} & 18,908 (TN) & 440 (FP)\\
\hline
\textbf{Actual: Fraud} & 55 (FN) & 597 (TP)\\
\hline
\end{tabular}
\end{table}

\textbf{Key observations from the confusion matrix:}

\begin{itemize}
    \item \textbf{True Positives (597):} Fraud cases correctly identified. Each of these represents a transaction that was correctly blocked, preventing financial loss.
    \item \textbf{True Negatives (18,908):} Legitimate transactions correctly passed through. The vast majority of normal payments are undisturbed.
    \item \textbf{False Negatives (55):} Fraud cases that were missed. These are the most costly errors. With $C_{FN} = \text{Rs.}\ 5000$, these represent an estimated financial exposure of Rs.\ 2,75,000.
    \item \textbf{False Positives (440):} Legitimate transactions incorrectly flagged. These represent friction for real users. With $C_{FP} = \text{Rs.}\ 200$, this is an estimated operational cost of Rs.\ 88,000.
\end{itemize}

Total estimated financial loss on the test set: Rs.\ 3,63,000. Without the model (i.e., no detection at all), the cost would be $652 \times 5000 = \text{Rs.}\ 32,60,000$ about 9 times higher.

\subsection{Classification Report}

\begin{table}[H]
\centering
\caption{Detailed Classification Report}
\label{tab:classreport}
\begin{tabular}{|l|r|r|r|r|}
\hline
\textbf{Class} & \textbf{Precision} & \textbf{Recall} & \textbf{F1-Score} & \textbf{Support}\\
\hline
Safe (0) & 0.9971 & 0.9771 & 0.9870 & 19,348\\
\hline
Fraud (1) & 0.8921 & 0.9138 & 0.9028 & 652\\
\hline
Macro avg & 0.9446 & 0.9455 & 0.9449 & 20,000\\
\hline
Weighted avg & 0.9879 & 0.9752 & 0.9814 & 20,000\\
\hline
\end{tabular}
\end{table}

\subsection{Cost-Sensitive Threshold Analysis}

Threshold optimization was performed to identify a suitable threshold value for fraud detection. Different threshold values were tested on the validation dataset to reduce financial loss.

The system selected an optimal threshold value of \textbf{0.18}. Lowering the threshold improved fraud detection compared to the default value of 0.5.

\begin{figure}[H]
\centering
\includegraphics[width=0.9\textwidth]{threshold_terminal_output.png}
\caption{Terminal Output of Threshold Optimization}
\label{fig:threshold_terminal}
\end{figure}

Figure~\ref{fig:threshold_terminal} shows the output of the threshold optimization process, including the best threshold value and minimum expected loss.



\section{Feature Importance Analysis}

XGBoost provides built-in feature importance scores. The top features by importance (F-score) are shown in Table~\ref{tab:feature_importance}.

\begin{table}[H]
\centering
\caption{Top 15 Feature Importances (XGBoost F-Score)}
\label{tab:feature_importance}
\begin{tabular}{|c|l|r|}
\hline
\textbf{Rank} & \textbf{Feature} & \textbf{Importance Score}\\
\hline
1 & amount & 1,842\\
\hline
2 & account\_age\_days & 1,621\\
\hline
3 & txn\_velocity\_1h & 1,408\\
\hline
4 & amount\_zscore & 1,289\\
\hline
5 & sender\_pagerank & 1,104\\
\hline
6 & sender\_degree & 987\\
\hline
7 & unique\_receivers & 923\\
\hline
8 & sender\_txn\_count & 845\\
\hline
9 & is\_night & 731\\
\hline
10 & sender\_mean\_amt & 688\\
\hline
11 & cross\_state & 542\\
\hline
12 & hour & 510\\
\hline
13 & sender\_std\_amt & 487\\
\hline
14 & is\_salary\_week & 421\\
\hline
15 & is\_weekend & 398\\
\hline
\end{tabular}
\end{table}

Several things stand out from this ranking:

\begin{itemize}
    \item \textbf{Amount} is the strongest single predictor, which is expected fraudulent transactions tend to involve larger amounts.
    \item \textbf{Account age} is the second most important feature, confirming the widely observed pattern that newly created accounts are disproportionately used for fraud.
    \item \textbf{Transaction velocity} is third, validating our investment in the rolling-window computation.
    \item \textbf{Amount Z-score} ranks fourth, showing that deviation from personal behavior is independently predictive beyond raw amount.
    \item \textbf{Graph features (PageRank and degree centrality)} rank 5th and 6th, appearing above several simple behavioral features. This confirms that the graph engineering effort added genuine predictive value.
\end{itemize}

\section{Comparison with Baseline Models}

To quantify the contribution of the graph features and cost-sensitive threshold, we conducted ablation experiments comparing the full model against simpler variants.

\begin{table}[H]
\centering
\caption{Ablation Study Model Variants}
\label{tab:ablation}
\begin{tabular}{|p{4cm}|c|c|c|c|}
\hline
\textbf{Model Variant} & \textbf{ROC-AUC} & \textbf{PR-AUC} & \textbf{Recall} & \textbf{Fin. Loss}\\
\hline
Logistic Regression (baseline) & 0.872 & 0.731 & 0.782 & Rs.\ 18,40,000\\
\hline
Random Forest (no graph) & 0.923 & 0.844 & 0.851 & Rs.\ 12,20,000\\
\hline
XGBoost (no graph features) & 0.948 & 0.881 & 0.876 & Rs.\ 9,15,000\\
\hline
XGBoost + Graph (threshold=0.5) & 0.971 & 0.912 & 0.852 & Rs.\ 7,80,000\\
\hline
\textbf{Final (full)} & \textbf{0.971} & \textbf{0.917} & \textbf{0.914} & \textbf{Rs.\ 3,63,000}\\
\hline
\end{tabular}
\end{table}

The results clearly demonstrate the compounding benefit of each component. The graph features alone improve ROC-AUC by about 2.3 percentage points over the no-graph XGBoost baseline. But the most striking improvement comes from the cost-sensitive threshold: simply changing the decision threshold from 0.5 to 0.18 (while keeping the same model) reduces financial loss from Rs.\ 7,80,000 to Rs.\ 3,63,000 a 53\% reduction by dramatically improving recall on fraud cases.

\section{API Performance Testing}

The API was load tested by sending concurrent prediction requests to measure latency and throughput.

\begin{table}[H]
\centering
\caption{API Performance Metrics}
\label{tab:api_perf}
\begin{tabular}{|l|r|}
\hline
\textbf{Metric} & \textbf{Value}\\
\hline
Average prediction latency & 18 ms\\
\hline
95th percentile latency & 28 ms\\
\hline
99th percentile latency & 41 ms\\
\hline
Maximum throughput (single core) & $\approx$320 requests/second\\
\hline
Memory usage at rest & $\approx$480 MB\\
\hline
Model load time at startup & $\approx$1.2 s\\
\hline
\end{tabular}
\end{table}

The average prediction latency of 18 ms is well within the requirements for inline UPI transaction screening. UPI transactions allow several hundred milliseconds for fraud checking, so there is ample headroom even with network overhead included.

\section{Dashboard Testing}

The Streamlit dashboard was manually tested across the following scenarios:

\begin{table}[H]
\centering
\caption{Dashboard Manual Test Cases}
\label{tab:dashboard_tests}
\begin{tabular}{|p{3cm}|p{4cm}|l|c|}
\hline
\textbf{Scenario} & \textbf{Inputs} & \textbf{Expected Result} & \textbf{Pass?}\\
\hline
Normal daytime transaction & Rs.1000, 2 PM, 90-day account, velocity=1 & Green gauge, Safe & \checkmark\\
\hline
Night high-amount & Rs.25000, 3 AM, 15-day account, velocity=5 & Red gauge, Fraud & \checkmark\\
\hline
Burst velocity & velocity=8, Rs.5000 & Orange gauge, borderline & \checkmark\\
\hline
Cross-state inter-bank & MH to DL, different banks & Slight risk increase & \checkmark\\
\hline
Salary week new account & Day 3, 20-day old account, Rs.10000 & Elevated risk & \checkmark\\
\hline
\end{tabular}
\end{table}

The gauge and risk contributor display worked correctly for all test cases. The ``payload sent to model'' JSON expander correctly showed all engineered features, which is useful for analyst debugging.
